%%%%%%%%%%%%%%%%%%%%%%%%%%%%%%%%%%%%
% Tables
%%%%%%%%%%%%%%%%%%%%%%%%%%%%%%%%%%%%
% =========================
% Sum stat tables
% =========================

\begin{landscape}
\begin{table}[H]
\centering
\caption{CPS Summary Statistics \label{tab:sumstat1}}
\centering
\begin{threeparttable}
\begin{tabular}[t]{lcccc}
\toprule
\multicolumn{1}{c}{ } & \multicolumn{1}{c}{\textbf{Overall}} & \multicolumn{3}{c}{\textbf{By Generation}} \\
\cmidrule(l{3pt}r{3pt}){2-2} \cmidrule(l{3pt}r{3pt}){3-5}
\textbf{Characteristic} & \makecell[c]{\textbf{All Sample} \\N = 318,404} & \makecell[c]{\textbf{First} \\N=40,033} & \makecell[c]{\textbf{Second} \\N=199,294} & \makecell[c]{\textbf{Third} \\N=79,077}\\
\midrule
Female & 0.49 & 0.53 & 0.49 & 0.49\\
Asian & 0.65 & 0.96 & 0.73 & 0.31\\
Age & 8.4 (5.1) & 10.9 (4.5) & 8.3 (5.1) & 7.7 (5.0)\\
College Graduate:\ \ 	 Father & 0.52 & 0.59 & 0.52 & 0.50\\
College Graduate:\ \ 	 Mother & 0.52 & 0.56 & 0.51 & 0.52\\
Total Family Income\ \ 	 (1999 dollars) & 87,031 (84,797) & 75,815 (74,489) & 88,295 (88,411) & 89,436 (80,051)\\
\bottomrule
\end{tabular}
\begin{tablenotes}
\item[1] The samples include children ages 17 and below who live in intact families. First-generation Asian immigrant children that were born in a Asian country. Native-born second-generation Asian immigrant children with at least one parent born in a Asian country. Finally, native-born third generation Asian immigrant children with native-born parents and at least one grand parent born in a Asian country.
\item[2] Data source is the 2004-2021 Current Population Survey.
\end{tablenotes}
\end{threeparttable}
\end{table}

\end{landscape}

\newpage
\pagebreak

\begin{table}

\caption{Current Population Survey (CPS) Summary Statistics \label{tab:sumstat-adults}}
\centering
\begin{threeparttable}
\begin{tabular}[t]{lccc}
\toprule
\multicolumn{1}{c}{ } & \multicolumn{1}{c}{\textbf{Overall}} & \multicolumn{2}{c}{\textbf{By Generation}} \\
\cmidrule(l{3pt}r{3pt}){2-2} \cmidrule(l{3pt}r{3pt}){3-4}
\textbf{Characteristic} & \makecell[c]{\textbf{All Sample} \\N = 48,153} & \makecell[c]{\textbf{First} \\N=35,728} & \makecell[c]{\textbf{Second} \\N=12,425}\\
\midrule
Female & 0.56 & 0.58 & 0.50\\
Asian & 0.90 & 0.96 & 0.73\\
Age & 46 (17) & 49 (16) & 38 (17)\\
Years of Education & 14.0 (3.5) & 13.9 (3.8) & 14.5 (2.5)\\
Total Family Income\ \ 	 (1999 dollars) & 68,121 (71,266) & 66,093 (70,444) & 73,950 (73,269)\\
\bottomrule
\end{tabular}
\begin{tablenotes}
\item[1] The samples include people of Asian ancestry ages 18 and above. First-generation Asian immigrants were born in a Asian country. Native-born second-generation Asian immigrants have at least one parent born in a Asian country. 
\item[2] Data source is the 2004-2021 Current Population Survey.
\end{tablenotes}
\end{threeparttable}
\end{table}


\newpage
\pagebreak
% =========================
% Identity by proxy
% =========================

\begin{table}[H]
\centering\centering
\caption{Main Effect of Proxy on Second-Generation's Asian Self-identification \label{tab:hispbyproxy}}
\centering
\fontsize{12}{14}\selectfont
\begin{threeparttable}
\begin{tabular}[c]{>{}lllll}
\toprule
Parents Type & All & Asian-Asian & Asian-White & White-Asian\\
\midrule
\textbf{Proxy:} &  &  &  & \\
\hspace{1em}\textbf{Mother} & 0.72 & 0.97 & 0.37 & 0.3\\
\hspace{1em}\textbf{Father} & 0.72 & 0.97 & 0.39 & 0.29\\
\hspace{1em}\textbf{Self} & 0.87 & 0.97 & 0.23 & 0.31\\
\hspace{1em}\textbf{Others} & 0.88 & 0.96 & 0.6 & 0.54\\
\bottomrule
\end{tabular}
\begin{tablenotes}
\item[1] The samples include children ages 17 and below who live in intact families. A proxy is the person that answered the Current Population Survey questionnaire.
\item[2] Data source is the 1994-2021 Current Population Survey.
\end{tablenotes}
\end{threeparttable}
\end{table}


\pagebreak
\newpage

% =========================
% Number of observations by generation
% =========================

\begin{table}[H]
\centering
\small\setstretch{1.0}
\caption{Self-Reported Racial Identity of Children of Asian Ancestry by Generation \label{tab:asianbygen}}
\begin{threeparttable}
\begin{tabular}{lccccc}
\toprule
 & & \multicolumn{4}{c}{Share self-reporting each identity} \\
\cmidrule(l{3pt}r{3pt}){3-6}
 & Observations & Asian only & White only & \specialcell{Asian and \\ White} & Other\\
\midrule
\textbf{First generation} & 15,499 & 0.94 & 0.03 & 0.01 & 0.02\\
\textbf{Second generation} & 80,137 & 0.72 & 0.10 & 0.14 & 0.04\\
\hspace{1em}Both parents Asian-born & 50,835 & 0.96 & 0.02 & 0.01 & 0.02\\
\hspace{1em}One Asian-born parent & 29,302 & 0.31 & 0.25 & 0.37 & 0.07\\
\hspace{1em}\hspace{1em}Asian father--White mother & 9,055 & 0.36 & 0.27 & 0.29 & 0.07\\
\hspace{1em}\hspace{1em}White father--Asian mother & 20,247 & 0.29 & 0.24 & 0.40 & 0.07\\
\textbf{Third generation} & 33,442 & 0.29 & 0.35 & 0.27 & 0.09\\
\hspace{1em}Four Asian grandparents & 5,744 & 0.93 & 0.04 & 0.01 & 0.02\\
\hspace{1em}Three Asian grandparents & 567 & 0.74 & 0.10 & 0.13 & 0.03\\
\hspace{1em}Two Asian grandparents & 12,678 & 0.24 & 0.28 & 0.39 & 0.09\\
\hspace{1em}One Asian grandparent & 14,453 & 0.07 & 0.54 & 0.26 & 0.12\\
\bottomrule
\end{tabular}
\begin{tablenotes}
\small
\item[1] The samples include children ages 17 and below who live in intact families. First-generation Asian immigrant children were born in an Asian country. Native-born second-generation children have at least one parent born in an Asian country; ``One Asian-born parent'' rows require the other parent to be born outside Asia (the AW and WA labels follow the paper's convention). Native-born third-generation children have native-born parents and at least one grandparent born in an Asian country, grouped by the number of Asian-born grandparents.
\item[2] Columns report the share of each row's observations in the racial identity categories used in the multinomial logit: ``Asian only,'' ``White only,'' and ``Asian and White'' (any combination of Asian and White, with or without Pacific Islander). ``Other'' collects Asian or Pacific Islander, other single races, and multiracial combinations without both an Asian and a White component; the multinomial logit excludes this category. Shares sum to one across the four columns.
\item[3] Data source is the 2004--2021 Current Population Survey.
\end{tablenotes}
\end{threeparttable}
\end{table}


\newpage
\pagebreak

\begin{table}[H]
\centering
\small\setstretch{1.0}
\caption{Self-Reported Racial Identity of Adults of Asian Ancestry by Generation \label{tab:asianbygen-adults}}
\begin{threeparttable}
\begin{tabular}{lccccc}
\toprule
 & & \multicolumn{4}{c}{Share self-reporting each identity} \\
\cmidrule(l{3pt}r{3pt}){3-6}
 & Observations & Asian only & White only & \specialcell{Asian and \\ White} & Other\\
\midrule
\textbf{First generation} & 286,553 & 0.95 & 0.03 & 0.00 & 0.02\\
\textbf{Second generation} & 102,372 & 0.70 & 0.15 & 0.10 & 0.05\\
\hspace{1em}Both parents Asian-born & 62,815 & 0.93 & 0.03 & 0.01 & 0.03\\
\hspace{1em}One Asian-born parent & 39,557 & 0.33 & 0.34 & 0.24 & 0.09\\
\hspace{1em}\hspace{1em}Asian father--White mother & 14,289 & 0.36 & 0.36 & 0.18 & 0.10\\
\hspace{1em}\hspace{1em}White father--Asian mother & 25,268 & 0.32 & 0.33 & 0.27 & 0.09\\
\bottomrule
\end{tabular}
\begin{tablenotes}
\small
\item[1] The samples include people of Asian ancestry ages 18 and above. First-generation Asian immigrants were born in an Asian country. Native-born second-generation Asian Americans have at least one parent born in an Asian country; ``One Asian-born parent'' rows require the other parent to be born outside Asia (the AW and WA labels follow the paper's convention). Grandparents' birthplaces are not observed for adults, so the third generation cannot be identified.
\item[2] Columns report the share of each row's observations in the racial identity categories used in the multinomial logit: ``Asian only,'' ``White only,'' and ``Asian and White'' (any combination of Asian and White, with or without Pacific Islander). ``Other'' collects Asian or Pacific Islander, other single races, and multiracial combinations without both an Asian and a White component; the multinomial logit excludes this category. Shares sum to one across the four columns.
\item[3] Data source is the 2004--2021 Current Population Survey.
\end{tablenotes}
\end{threeparttable}
\end{table}


\pagebreak
\newpage


% =========================
% Migration table
% =========================

\begin{table}[H]
\centering\centering
\caption{Relationship Between Bias and Migration \label{regtab-mig-01}}
\centering
\resizebox{0.7\textwidth}{!}{
\begin{threeparttable}
\begin{tabular}[t]{lccc}
\toprule
  & \specialcell{(1) \\ Migrated from \\ Birth Place} & \specialcell{(2) \\ Migrated from \\ Birth Place} & \specialcell{(3) \\ $Bias_{ist} - Bias_{ilb}$}\\
\midrule
$Bias_{st}$ & 0.13* &  & \\
 & (0.07) &  & \\
$Bias_{lb}$ &  & -0.03 & \\
 &  & (0.17) & \\
Asian &  &  & 0.02\\
 &  &  & (0.04)\\
Female & 0.00 & -0.01 & 0.00\\
 & (0.00) & (0.00) & (0.02)\\
College Graduate: Mother & 0.01*** & 0.00 & -0.01\\
 & (0.00) & (0.01) & (0.03)\\
College Graduate: Father & -0.03*** & -0.03*** & 0.03\\
 & (0.01) & (0.01) & (0.02)\\
\midrule
Observations & 73,563 & 41,641 & 2,075\\
Mean & 0.15 & 0.15 & -0.1\\
Year $\times$ Region FE & X &  & \\
Birthyear $\times$ Birth Region FE &  & X & \\
\bottomrule
\multicolumn{4}{l}{\rule{0pt}{1em}* p $<$ 0.1, ** p $<$ 0.05, *** p $<$ 0.01}\\
\end{tabular}
\begin{tablenotes}
\small
\item[1] \footnotesize{Columns (1) and (2) estimate: $BirthPlaceMigration_{i} = \beta_1 AntiAsianBias + X_{i}\pi + \gamma + \varepsilon_{i}$, where $BirthPlaceMigration$ indicates whether person $i$ lives in a state different from their birth state. Column (1) uses bias in the current state at survey time; column (2) uses bias in the birth state at birth year. Column (3) estimates: $Y_{i} = \beta_0 + \beta_1 Asian_{i} + X_{i}\pi + \varepsilon_{i}$, where $Y_{i} \equiv AntiAsianBias_{ist} - AntiAsianBias_{ilb}$ is the difference between bias in the current state and bias in the birth state, and $Asian_{i}$ is self-reported Asian identity. I include controls for sex, quartic age, parental education, fraction of Asians in a state, and region $\times$ year fixed effects (column 1) or birthyear $\times$ birth region fixed effects (column 2). Standard errors are clustered on the state level.}
\item[2] \footnotesize{The samples include children ages 17 and below who live in intact families. 
                      Native-born second-generation Asian immigrant children with both
                      parents born in an Asian country. The sample in the column (3) regression is further restricted to only those that migrated from their birth state.}
\item[3] \footnotesize{Data source is the 2004--2021 American Community Survey (IPUMS USA).}
\end{tablenotes}
\end{threeparttable}}
\end{table}


\pagebreak
\newpage

% =========================
% Marriage table
% =========================


\begin{table}[H]
\centering\centering
\caption{Relationship Between Bias and Interracial Marriages \label{regtab-logit-02}}
\centering
\begin{threeparttable}
\begin{tabular}[t]{lccc}
\toprule
\multicolumn{2}{c}{ } & \multicolumn{1}{c}{Asian Men} & \multicolumn{1}{c}{Asian Women} \\
\cmidrule(l{3pt}r{3pt}){3-3} \cmidrule(l{3pt}r{3pt}){4-4}
  & \specialcell{(1) \\ Interracial} & \specialcell{(2) \\ Interracial} & \specialcell{(3) \\ Interracial}\\
\midrule
Bias & $0.04$ & $-0.01$ & $0.03$\\
 & ($0.03$) & ($0.03$) & ($0.03$)\\
College Graduate: Wife & $0.04$*** & $0.04$*** & $0.05$***\\
 & ($0.01$) & ($0.01$) & \vphantom{1} ($0.01$)\\
College Graduate: Husband & $-0.01$ & $-0.01$ & $-0.02$*\\
 & ($0.01$) & ($0.01$) & ($0.01$)\\
\midrule
Observations & $69,800$ & $52,103$ & $60,214$\\
Year $\times$ Region FE & X & X & X\\
\bottomrule
\multicolumn{4}{l}{\rule{0pt}{1em}* p $<$ 0.1, ** p $<$ 0.05, *** p $<$ 0.01}\\
\end{tabular}
\begin{tablenotes}
\small
\item[1] \footnotesize{Each column estimates the following linear probability model: $interracial_{ist} = \beta_1 AntiAsianBias_{st} + X_{ist}\pi + \gamma_{rt} + \varepsilon_{ist}$, where $interracial_{ist}$ indicates a couple in which one partner was born in an Asian country and the other was not (labelled Asian husband--White wife or White husband--Asian wife, following the paper's convention), $AntiAsianBias_{st}$ is average anti-Asian bias in state $s$ at time $t$, and $X_{ist}$ includes partner-specific controls (sex, age, education, and years since immigrating to the United States). The unit of observation is a second-generation child; columns (2) and (3) restrict to children of Asian-born fathers and Asian-born mothers, respectively. Standard errors are clustered on the state level.}
\item[2] \footnotesize{Data source is the 2004-2020 Current Population Survey Data.}
\end{tablenotes}
\end{threeparttable}
\end{table}


\pagebreak
\newpage

% =========================
% Heterogeneity table
% (replaces the interaction
% coefficient figures)
% =========================

\begin{table}[H]
\centering
\small\setstretch{1.0}
\caption{Heterogeneity in the Effect of Anti-Asian Bias on Asian Identification \label{tab:heterogeneity-interactions}}
\resizebox{\ifdim\width>\linewidth\linewidth\else\width\fi}{!}{
\begin{threeparttable}
\begin{tabular}{lcccc}
\toprule
\multicolumn{5}{l}{\textit{Panel A: Mixed-race second generation (children)}} \\
 & Bias $\times$ Female & Bias $\times$ Mother college & Bias $\times$ Father college & Observations \\
\midrule
Asian father--White mother & $-0.005$ & $-0.225$** & $0.039$ & $9{,}055$ \\
 & ($0.052$) & ($0.097$) & ($0.094$) & \\
White father--Asian mother & $-0.011$ & $0.006$ & $-0.036$ & $20{,}247$ \\
 & ($0.048$) & ($0.047$) & ($0.035$) & \\
\textbf{Pooled} & $-0.012$ & $-0.067$* & $-0.024$ & $29{,}302$ \\
 & ($0.052$) & ($0.034$) & ($0.037$) & \\
\midrule
\multicolumn{5}{l}{\textit{Panel B: Mixed-race second generation (adults)}} \\
 & Bias $\times$ Female & Bias $\times$ College graduate & Bias $\times$ Above-median income & Observations \\
\midrule
Asian father--White mother & $-0.026$ & $-0.084$* & $-0.006$ & $1{,}710$ \\
 & ($0.047$) & ($0.047$) & ($0.048$) & \\
White father--Asian mother & $0.017$ & $0.003$ & $-0.022$ & $2{,}976$ \\
 & ($0.042$) & ($0.041$) & ($0.041$) & \\
\textbf{Pooled} & $-0.012$ & $-0.023$ & $-0.020$ & $4{,}686$ \\
 & ($0.033$) & ($0.039$) & ($0.040$) & \\
\midrule
\multicolumn{5}{l}{\textit{Panel C: Third generation (children)}} \\
 & Bias $\times$ Female & Bias $\times$ Mother college & Bias $\times$ Father college & Observations \\
\midrule
One Asian grandparent & $-0.010$ & $-0.023$ & $-0.034$ & $14{,}453$ \\
 & ($0.028$) & ($0.025$) & ($0.030$) & \\
Two Asian grandparents & $0.034$ & $-0.077$ & $0.000$ & $12{,}678$ \\
 & ($0.028$) & ($0.061$) & ($0.073$) & \\
Three Asian grandparents & $0.019$ & $-0.297$ & $0.096$ & $567$ \\
 & ($0.077$) & ($0.211$) & ($0.143$) & \\
\textbf{Pooled} & $-0.004$ & $-0.067$** & $-0.005$ & $27{,}698$ \\
 & ($0.017$) & ($0.030$) & ($0.039$) & \\
\bottomrule
\end{tabular}
\begin{tablenotes}
\small
\item[1] \footnotesize{Each row is a separate weighted linear probability model, $A_{ist} = \beta_1 Bias_{st} + \sum_k \delta_k (Bias_{st} \times Z_{ist}^k) + \sum_k \theta_k Z_{ist}^k + \mathbf{X}_{ist}\pi + \gamma_{rt} + \varepsilon_{ist}$, where $A_{ist}$ is self-reported Asian identity, $Bias_{st}$ is the demeaned state-level anti-Asian bias index (standard deviation units), and the table reports the interaction coefficients $\delta_k$. In Panels A and C, $Z$ is female, mother college graduate, and father college graduate; in Panel B, $Z$ is female, own college degree (16 or more years of schooling), and family income above the median of the full adult file. All models include region $\times$ year fixed effects and controls for a quartic in age and the state Asian population share, and are weighted by CPS person weights. Panel A uses children ages 17 and younger living with both parents, with one parent born in an Asian country and the other born outside Asia (labelled White, following the paper's convention; about 86 percent of these parents are US-born); Panel B uses adults ages 18 and older with the same parental configuration, restricted to March (ASEC) respondents because family income is observed only for them; Panel C uses native-born children with native-born parents and one, two, or three grandparents born in an Asian country. The pooled rows do not include group indicators, matching the specifications previously plotted. Standard errors clustered at the state level in parentheses. $^{*}p<0.1$, $^{**}p<0.05$, $^{***}p<0.01$. Data source is the 2004--2021 Current Population Survey.}
\end{tablenotes}
\end{threeparttable}}
\end{table}


\pagebreak
\newpage


% =========================
% Diagram of how
% generations are defined
% =========================

\pagebreak
\newpage

\begin{center}
\begin{figure}[H]
\caption{Diagram of the Three Different Generations of Asian Immigrants.}
\includegraphics[width=\textwidth, height=9cm]{diag.png} 
\label{fig:diag}
\end{figure}
\hfill%
\end{center}

% =========================
% Bias and Self-reported 
% Asian Identity in the Least 
% and Most Biased Places
% =========================


\begin{center}
\begin{figure}[H]
\caption{Bias and Self-reported Asian Identity in the Least and Most Biased Places} \label{fig:toptwobias}

% first
\begin{subfigure}{.9\textwidth}
\caption{Bias Over Time}
\centering
\includegraphics[width=.9\linewidth]{Bias_twostates.png} 
\label{fig:skiniat}
\end{subfigure}
% Second
\begin{subfigure}{.9\textwidth}
\caption{Self-reported Asian Identity Over Time}
\centering
\includegraphics[width=.9\linewidth]{Bias_twostates-asian.png} 
\label{fig:Asian-twostates}
\end{subfigure}
\caption*{\footnotesize{These two panels show the trends in bias (panel a) and self-reported Asian identity among Asian immigrants (panel b) of the least and most biased places in the data. The District of Colombia is the least biased geographical area, and North Dakota is the most biased. The bias units are in standard deviations. Self-reported Asian identity is among first, second, and third-generation Asian immigrants aged 17 and younger still living in intact families.\\
Bias data is from the 2004--2021 Harvard's Project Implicit Association Test scores, American National Election Studies (ANES), and state-level hate crimes against Asians. Identity data is from the 2004--2021 Current Population Survey (CPS).}}
\end{figure}
\end{center}


\pagebreak
\newpage

% =========================
% Multinomial Logit Marginal
% Effects: Pooled Samples
% =========================

\begin{figure}[H]
\centering
\caption{Marginal Effects of Key Covariates on Racial Identity Choice}
\label{fig:marginal-effects-pooled}

\begin{subfigure}{.76\textwidth}
\caption{Mixed-Race Second-Generation Asian Americans}\label{fig:marginal-effects-second-pooled}
\centering
\includegraphics[width=1\linewidth]{boot_me_second_AW_WA.png}
\end{subfigure}

\vspace{0.3cm}

\begin{subfigure}{.76\textwidth}
\caption{Third-Generation Asian Americans}\label{fig:marginal-effects-third-pooled}
\centering
\includegraphics[width=1\linewidth]{boot_me_third_oneTwoThreeAsian.png}
\end{subfigure}

\caption*{\footnotesize{This figure displays average marginal effects of key covariates on the probability of each racial identity outcome, estimated from equation (\ref{eq:multinomial_logit}). Panel (a) pools second-generation children with one parent born in an Asian country and one parent born outside Asia; panel (b) pools third-generation children with one, two, or three Asian-born grandparents. For anti-Asian bias, the marginal effect is the derivative of each predicted probability with respect to the bias index, evaluated for every child at that child's own covariate values and averaged over the estimation sample, so it is the average change in probability per one-unit (one standard deviation) increase in bias. For the binary covariates, it is the average change in predicted probability from switching the indicator from 0 to 1 for every child. Models include region $\times$ year fixed effects and controls for a quartic in age and the state Asian population share, and are weighted by CPS person weights. Each point is the median of the average marginal effect across 1{,}000 bootstrap resamples of children, drawn with replacement within region $\times$ year cells; the bars span the 2.5th and 97.5th percentiles of the same resamples. Because children rather than states are resampled, the intervals do not account for dependence within state-year cells.}}
\end{figure}
